\documentclass[../main.tex]{subfiles}

\begin{document}

\section{Introduction}
\label{sec:intro}

\subsection{Purpose}

\subsection{Scope}

Out of scope:
\begin{itemize}
    \item PHDF5
    \item HDFView
    \item HSDS
    \item AI
\end{itemize}

\subsubsection{Safety}

\subsubsection{Security}

\subsubsection{Privacy}
The National Institute of Standards and Technology (NIST) defines privacy as “a condition that safeguards human autonomy and dignity through various means, including confidentiality, predictability, manageability, and disassociability.” Historically, privacy was interpreted as preventing unauthorized access to or disclosure of Personally Identifiable Information (PII). However, NIST emphasizes that contemporary privacy risks extend far beyond traditional notions of confidentiality. Modern privacy concerns increasingly arise from inference, aggregation, behavioral profiling, secondary use of data, context changes, and surveillance effects [1]. NIST further emphasizes that privacy risks are distinct from cybersecurity risks and may emerge even when systems are secure, lawful, and functioning exactly as intended [2]. HDF5 inherits many of the same classes of privacy risks identified by NIST, although these risks manifest differently in scientific and industrial environments. 
HDF5 was designed to maximize data shareability, extensibility, portability, and long-term usability and not to enforce privacy data protections. Its rich metadata model, self-describing structure, plugin architecture, and interoperability layers can expose information extending beyond the raw data stored within the file. Sensitive information in HDF5 ecosystems may emerge through HDF5 internal and user-defined metadata, for example, HDF5 attributes, object names, file organization, and objects properties, as well as through HDF5 functionality, such as different types of connectors and plugins. It may be exposed during HDF5 data lifecycle through provenance records, workflow artifacts, and correlations across multiple datasets and files. 
A common misconception in scientific computing is that privacy concerns are resolved once data access is restricted or datasets are anonymized. However, privacy risks often persist even when traditional security mechanisms are correctly implemented. Restricting access does not prevent inference from metadata. Anonymization may remove direct identifiers while leaving contextual information sufficient for re-identification. Similarly, data reduction or aggregation does not eliminate the possibility that sensitive information may be reconstructed through correlations across files, workflows, or external repositories. In modern scientific ecosystems, privacy exposure frequently results from accumulation of data and metadata over time rather than disclosure of a single sensitive piece of raw data or metadata.
Privacy risks are further amplified by the scientific community's increasing emphasis on data openness and sharing. HDF5 datasets routinely move across organizations, computing systems, cloud platforms, workflows, and both public and private repositories. Data are frequently repurposed beyond their original intent and analyzed by users who were not involved in their creation. Machine learning systems further amplify these concerns by discovering previously unknown relationships within data stored in HDF5 and enabling new forms of inference that may not have been anticipated when the data were originally collected. Consequently, the privacy characteristics of data stored in HDF5 files may evolve throughout the data lifecycle.
Therefore, privacy exposure in HDF5 is not solely a matter of unauthorized access, data breaches, misuse, or malicious behavior; it is fundamentally different from traditional notions of data security and safety. Neither the HDF5 file format nor the HDF5 software stack guarantees privacy protection, and both may present privacy risks. Moreover, privacy risks may emerge during legitimate scientific workflows through data sharing, reuse, integration, and analysis. In this document, we examine and classify the various privacy exposure mechanisms, or exposure surfaces, present in HDF5 with the goal of increasing awareness of privacy risks associated with the use of HDF5.


\section{Motivation}
\label{sec:motivation}

\begin{figure}[h!]
    \centering
    \includegraphics[width=\textwidth]{images/HDF5_Surfaces.png}
    \caption{HDF5 vulnerability sites (interfaces: rounded rectangles)}
\end{figure}

\subsection{Problem Statement}

\subsection{Goals}

\begin{itemize}
    \item Emphasize broader applicability for comparable OSEs
\end{itemize}

\end{document}
